\chapter{Conclusion}

Based on the evidence, discussed in previous chapters, it is reasonable to presume that given the high engagements of developers, their low comprehension confidence, and
a clear preference for \textit{accessible} sources, the conclusion can be made that no single existing academic medium, whether formal or informal, can effectively serve
the needs of the game development community.

\section{Is there a requirement for a new model of paper?}

Whilst the current academic ecosystem seems to be functioning as is, the argument could, and perhaps \textit{should} be made whether a new 'model' of paper standard should be
proposed. This new model could provide a hybrid approach where the rigour, structure, and formality of peer-reviewed papers is married to the clarity and practical approach of
blogs, or other less formal types of academia.

\section{What would this paper look like?}

This new model's function should aim to actively \textit{teach} its concepts using different learning techniques, and should perhaps contain:

\begin{itemize}
    \item \textbf{Theoretical Context:} A concise summary of the core mathematics
    \item \textbf{Visualisation:} Step-by-step (\textit{or at the least, more clear and informative}) diagrams that help map abstract concepts into more tangible, real-world examples
    \item \textbf{Implementations:} Functional, and annotated pseudocode that help translate the theory and concepts discussed
\end{itemize}

The necessity of a new model can be argued. It should be noted that it is unreasonable to expect all academia, to adopt a new standard of paper,
however papers and specific fields that have a large game-developer audience may benefit.